\section{The derivative as a function}

The derivative can be computed at a single point. But often we want to know the derivative at many points. When we assign to each input \(x\) the derivative of \(f\) at \(x\), we obtain a new function.

This new function is called the derivative function and is written as
\[
f'(x).
\]

\subsection*{From \(f'(a)\) to \(f'(x)\)}

The number \(f'(a)\) is the derivative at one point. It measures the local rate of change at \(x=a\). The expression \(f'(x)\) describes the derivative at a variable point \(x\). It tells us how the local rate of change depends on the input.

\begin{examplebox}
For
\[
f(x)=x^2,
\]
we found from the definition that
\[
f'(a)=2a.
\]
Since the point \(a\) was arbitrary, we can write
\[
f'(x)=2x.
\]
This means that the slope of the tangent line to \(y=x^2\) at the point with input \(x\) is \(2x\).
\end{examplebox}

The derivative function is therefore a function about another function. It describes how the original function changes.

\subsection*{Different notations}

Several notations are commonly used for derivatives. If \(y=f(x)\), then we may write
\[
f'(x),
\qquad
 y',
\qquad
\frac{dy}{dx},
\qquad
\frac{d}{dx}f(x).
\]
These notations emphasize different aspects of the same idea.

The notation \(f'(x)\) emphasizes that the derivative is a function. The notation \(dy/dx\) emphasizes a ratio of changes in output and input. The notation \(\frac{d}{dx}\) emphasizes the operation of differentiating with respect to \(x\).

\begin{remarkbox}
The notation changes, but the mathematical object is the same: the derivative describes local rate of change.
\end{remarkbox}

\subsection*{Where the derivative is defined}

The derivative function may not be defined everywhere that the original function is defined. A function may have corners, cusps, vertical tangents, or discontinuities where the derivative does not exist.

\begin{examplebox}
The function
\[
f(x)=|x|
\]
is defined for all real numbers. But its derivative does not exist at \(x=0\), because the left-hand slope is \(-1\) and the right-hand slope is \(1\).

For \(x>0\), the function is \(f(x)=x\), so \(f'(x)=1\). For \(x<0\), the function is \(f(x)=-x\), so \(f'(x)=-1\). At \(x=0\), there is no derivative.
\end{examplebox}

Thus the domain of \(f'\) can be smaller than the domain of \(f\).

\subsection*{Differentiability and continuity}

If a function is differentiable at a point, then it is continuous at that point. This means that having a derivative is a stronger condition than being continuous.

\begin{theorembox}
If \(f\) is differentiable at \(a\), then \(f\) is continuous at \(a\).
\end{theorembox}

The converse is not true. A function may be continuous but not differentiable. The function \(f(x)=|x|\) is continuous at \(0\), but it is not differentiable there.

\begin{summarybox}
When reading a derivative function, ask:
\begin{itemize}
\item Is the derivative being evaluated at one point or treated as a function?
\item What notation is being used for the derivative?
\item On which inputs does the derivative exist?
\item Does the graph have corners or breaks where differentiability may fail?
\item What does the derivative say about local change?
\end{itemize}
\end{summarybox}

The derivative function is a tool for reading how the original function changes from point to point.